\documentclass[a4paper,11pt]{article}
\usepackage{exam-style-341}
\examinehead{341 农业综合知识三 · 网络技术分章练习}{第09章 网络安全 参考答案与解析}
\begin{document}

\unitanswerscover{341 农业综合知识三 · 网络技术分章练习}{第09章 网络安全 参考答案}{第09章练习的答案与解析；判断题均给出改正后的表述，RSA 计算给出完整推导与验证过程，综合题按“威胁—防护—溯源”三层给出要点。}

\answerpart{一、单项选择题（10 题，每题 3 分，共 30 分）}

\q C。\begin{analysis}被动攻击指攻击者只窃听、分析通信内容而不改变数据流，截获（窃听）是典型代表，还包括流量分析（教材 7.1.1，印刷页 334—335）。A 篡改、B 伪造、D 拒绝服务都会改变数据或破坏服务，属于主动攻击。\end{analysis}

\q A。\begin{analysis}DES（数据加密标准）加密与解密使用同一密钥，属于对称密钥密码体制，分组 64 位、有效密钥 56 位（教材 7.2.1，印刷页 338）。B 公钥密码体制的典型算法是 RSA；C 散列函数如 MD5、SHA-1；D 数字签名通常基于公钥密码实现，DES 本身不是签名算法。\end{analysis}

\q C。\begin{analysis}公钥密码体制中，发送方用接收方的公钥加密，只有接收方用自己的私钥才能解密，从而实现保密通信（教材 7.2.2，印刷页 339—340）。A 用 A 的公钥加密则只有 A 自己能解，不能发给 B；B、D 私钥不能公开用于加密（私钥用于解密或签名）。\end{analysis}

\q A。\begin{analysis}数字签名时，发送方用自己私钥对报文（或散列值）进行运算，接收方用发送方的公钥验证；由于只有发送方持有私钥，因此可鉴别发送者并防止抵赖（教材 7.3.1，印刷页 341）。B、C、D 的密钥使用关系均错误。\end{analysis}

\q B。\begin{analysis}MD5 是一种密码散列函数，输出 128 位散列值；SHA-1 输出 160 位（教材 7.3.1）。A 64 位是 DES 分组长度；C 160 位是 SHA-1 输出长度；D 256 位常见于 SHA-256 或 AES-256，均不是 MD5 的输出长度。\end{analysis}

\q B。\begin{analysis}SSL／TLS 位于应用层与运输层之间（通常归入运输层安全协议），为 HTTP 等应用提供加密、鉴别和完整性保护（教材 7.5.2，印刷页 357—358）。A 网络层的安全协议是 IPsec；C 会话层、D 应用层的安全协议（如 PGP、SET）与 SSL／TLS 不同。\end{analysis}

\q B。\begin{analysis}HTTPS 即 HTTP over TLS／SSL，默认端口为 443（教材 7.5.2）。A 80 是 HTTP 默认端口；C 8080 常用作 HTTP 代理或备用端口；D 21 是 FTP 控制连接端口。\end{analysis}

\q B。\begin{analysis}防火墙是部署在可信内部网络与不可信外部网络之间、按安全策略对进出通信实施访问控制的技术（教材 7.6.1，印刷页 364）。A 查杀病毒是防病毒软件的功能；C 加密由密码技术完成；D 邮件加密存储由应用层安全机制完成。\end{analysis}

\q C。\begin{analysis}拒绝服务攻击通过耗尽目标主机的带宽、连接数或处理能力，使其无法为合法用户提供服务，破坏的是可用性（教材 7.1.1，印刷页 335）。A 保密性对应窃听，B 完整性对应篡改，D 不可否认性对应数字签名，均不是拒绝服务攻击的主要目标。\end{analysis}

\q A。\begin{analysis}RSA 的安全性基于大整数分解的困难性：已知 n=pq 求 p、q 在计算上不可行，从而无法由公钥推出私钥（教材 7.2.2，印刷页 339）。B 离散对数问题是 ElGamal、ECC 等的基础；C、D 与 RSA 的安全性依据无关。\end{analysis}

\answerpart{二、多项选择题（3 题，每题 4 分，共 12 分）}

\q ABC。\begin{analysis}主动攻击会改变数据或伪造数据流，包括篡改、伪造、重放、拒绝服务等（教材 7.1.1）。D 错误：截获只窃听、不改变数据，属于被动攻击。\end{analysis}

\q ABC。\begin{analysis}DES、AES、IDEA 都是对称密钥密码算法。D 错误：RSA 是公开密钥密码算法，加密与解密使用不同密钥。\end{analysis}

\q ABC。\begin{analysis}数字签名应满足三项功能：报文鉴别（确认发送者）、报文完整性（报文未被篡改）、不可否认（发送方事后不能抵赖）（教材 7.3.1，印刷页 341—342）。D 错误：数字签名会增加计算与传输开销，不能提高传输速率。\end{analysis}

\answerpart{三、判断题（6 题，每题 2 分，共 12 分）}

\q √。\begin{analysis}对称密钥体制又称常规密钥体制，加密与解密使用同一密钥，因此密钥必须由通信双方共享并保密（教材 7.2.1）。\end{analysis}

\q √。\begin{analysis}公钥密码体制中，加密密钥（公钥）向公众公开，解密密钥（私钥）由接收方自己保密；由公钥在计算上不能推出私钥（教材 7.2.2）。\end{analysis}

\q ×。改正：数字签名使用发送方自己的私钥对报文（或其散列值）进行签名，接收方用发送方的公钥验证，而不是用接收方的公钥签名。\begin{analysis}若用接收方公钥签名，则任何人都能验证，无法确认发送者身份，也不能防止发送方抵赖。\end{analysis}

\q ×。改正：MD5 是一种密码散列函数（单向散列算法），不是对称加密算法；它用于产生报文摘要、检验完整性。\begin{analysis}散列函数不可逆，不能用于加密解密；MD5 输出 128 位散列值，常与数字签名配合使用。\end{analysis}

\q ×。改正：防火墙部署在内部网络与外部网络的边界上，主要防范来自外部网络的攻击，不能防止来自内部网络的攻击行为。\begin{analysis}防火墙是边界访问控制技术；内部人员的越权、泄露等威胁需要身份认证、审计、入侵检测等内部防护措施。\end{analysis}

\q √。\begin{analysis}报文鉴别包含两方面内容：鉴别报文发送者（实体鉴别）和鉴别报文完整性（报文未被篡改），可用 MAC 或数字签名实现（教材 7.3.1）。\end{analysis}

\answerpart{四、计算题（1 题，共 6 分）}

\q 参考答案与推导：
（1）模数 $n$ 与欧拉函数（2 分）：
\[
n=pq=3\times 11=33,\qquad \varphi(n)=(p-1)(q-1)=2\times 10=20
\]
（2）私钥 $d$（2 分）：求 $d$ 使 $e\cdot d\equiv 1\pmod{20}$，即 $3d\equiv 1\pmod{20}$。逐一试验：$3\times 7=21=20+1$，故 $d=7$。（也可用扩展欧几里得算法求得。）
（3）加密（1 分）：$C=M^{e}\bmod n=5^{3}\bmod 33=125\bmod 33$。因 $33\times 3=99$，$125-99=26$，故密文 $C=26$。
（4）解密验证（1 分）：
\[
C^{d}\bmod n=26^{7}\bmod 33
\]
先算 $26^{2}=676=33\times 20+16$，故 $26^{2}\equiv 16$；$26^{4}\equiv 16^{2}=256=33\times 7+25$，故 $26^{4}\equiv 25$；
$26^{7}=26^{4}\times 26^{2}\times 26\equiv 25\times 16\times 26=10400$，而 $33\times 315=10395$，$10400-10395=5$，故 $26^{7}\bmod 33=5=M$，解密还原明文 5。
\begin{analysis}本题只用于演示 RSA 的计算流程，实际使用的 p、q 是数百位的大素数。注意：$\varphi(n)$ 不是 $n-1$；私钥 d 与 e 关于 $\varphi(n)$ 互逆，而不是关于 n 互逆。\end{analysis}

\answerpart{五、简答题（2 题，每题 10 分，共 20 分）}

\q 参考答案要点：
（1）密钥使用（2 分）：对称密钥体制加密与解密使用同一密钥；公钥密码体制使用一对密钥，公钥公开、私钥保密，公钥加密私钥解密（或私钥签名公钥验证）。
（2）典型算法（2 分）：对称密钥体制有 DES、AES、IDEA 等；公钥密码体制有 RSA、ECC、ElGamal 等。
（3）优缺点（3 分）：对称密钥算法速度快、适合加密大量数据，但密钥分配与管理困难，n 个用户两两通信需要大量密钥；公钥算法无需事先共享密钥、便于密钥分配和数字签名，但运算速度慢、开销大。
（4）主要用途（3 分）：对称密钥用于数据加密（如文件、通信内容的批量加密）；公钥密码用于密钥交换（协商会话密钥）、数字签名与身份鉴别。实际系统常采用“公钥协商会话密钥＋对称密钥加密数据”的混合方式（教材 7.2）。

\q 参考答案要点：
（1）基本原理（3 分）：发送方用自己私钥对报文（或其散列值）进行运算得到签名，随报文一起发送；接收方用发送方公钥对签名进行验证运算并还原出散列值，与自己对报文计算的散列值比较，一致则签名有效（教材 7.3.1，印刷页 341）。
（2）三项功能（3 分）：报文鉴别（确认报文确由持有私钥的发送方发出）；报文完整性（报文未被篡改，篡改后散列值不匹配）；不可否认（只有发送方持有私钥，事后不能抵赖）。
（3）先 Hash 再签名（4 分）：公钥密码运算开销大，而报文可能很长，直接对整篇报文签名代价过高；散列函数把任意长度报文压缩成固定长度的短摘要（MD5 128 位、SHA-1 160 位），对摘要签名即可，签名短、效率高；同时散列函数的抗碰撞性保证报文改动后摘要必然改变，完整性仍可验证。因此实际中普遍采用“先 Hash、再对摘要签名”的方式。

\answerpart{六、综合题（1 题，共 20 分）}

\q 参考答案与答题框架：

\textbf{【破题句】}该系统的安全需求可归纳为“保密、完整、可用、可追溯、不可抵赖”，防护应按“威胁识别—协议与技术防护—溯源数据防篡改”三层设计。

\textbf{（1）主要安全威胁（6 分）}
\begin{analysis}
按两类给分，每类 3 分：
被动攻击（不改变数据，难以检测）：（1）窃听用户登录口令、身份信息与支付信息；（2）对网站与物联网设备的通信进行流量分析，获取经营与生产信息；（3）非法访问数据库读取溯源与客户数据（若未加密）。
主动攻击（改变或伪造数据/服务）：（1）篡改农产品溯源记录（产地、检测结果、流通环节）；（2）伪造身份进行网络钓鱼、冒充管理员或消费者；（3）重放攻击（重放支付请求或身份认证报文）；（4）拒绝服务攻击（攻击 Web 服务器或大量物联网设备，使其无法提供服务）；（5）恶意代码、勒索软件攻击服务器和终端。
\end{analysis}

\textbf{（2）Web 访问与在线支付的安全协议与技术（8 分）}
\begin{analysis}
每点 2 分，答出 4 点即可：
（1）采用 HTTPS（HTTP over SSL／TLS，端口 443）：对浏览器与服务器之间的通信加密，保证保密性与完整性；
（2）采用数字证书与 CA 认证：验证网站服务器身份，防止钓鱼网站冒充；
（3）采用强身份认证：口令加短信验证码／动态令牌（双因素认证），保护用户账户与支付操作；
（4）对支付请求、订单信息使用数字签名与完整性校验（MAC），防止篡改与抵赖；
（5）部署防火墙、入侵检测系统（IDS）与 Web 应用防火墙（WAF），过滤恶意流量；
（6）敏感数据（口令、支付信息、客户信息）加密存储，口令加盐散列保存，并定期备份。
\end{analysis}

\textbf{（3）溯源数据防篡改方案（6 分）}
\begin{analysis}
可采用“区块链＋数字签名＋散列链＋时间戳”的方案（答出要点即可）：
（1）各环节主体（农户、加工企业、物流、销售）用各自私钥对上传的溯源数据签名，接收方用其公钥验证，保证数据来源真实、不可抵赖；
（2）数据按时间顺序打包成区块，每个区块包含前一区块的散列值，形成哈希链；一旦修改某条记录，该区块及其后所有区块的散列值都会改变，容易被发现；
（3）区块由多个节点共同维护并经共识机制确认，单方难以篡改，实现分布式存储与不可篡改；
（4）配合可信时间戳与检测报告等第三方数据上链，消费者扫码即可查询全链条信息，实现“来源可查、去向可追、责任可究”。
\end{analysis}

\textbf{【给分点】}（1）被动攻击 3 分、主动攻击 3 分；（2）4 点各 2 分；（3）签名 2 分、散列链 2 分、共识与分布式 1 分、查询追溯 1 分。

\end{document}
